
In driving and other epoch-structured paradigms, it is common to apply a continuous-signal blink detector to the raw or minimally preprocessed recording and subsequently map the detected blinks back onto the epoch grid. Both BLINKER and the standard MNE pipeline follow this approach, in which channel-specific thresholds are derived from a signal pool containing both blink-heavy and blink-free epochs. The central premise of the proposed pipeline is that pooling these heterogeneous epochs dilutes the amplitude statistics used for threshold estimation, thereby distorting threshold placement and shifting the resulting precision--recall trade-off.
%MNE built-in plugin provides a second reference point grounded in community tooling, establishing its performance under an event-level blink criterion, for which it was not designed, and characterises the gap that specialised methods must close.
Accordingly, we hypothesise that naive concatenation will produce systematically lower session-level $F_1$ scores than the epoch-aware pipeline because of a less favourable precision--recall trade-off. To evaluate this hypothesis, all four pipelines (i.e., BLINKER, MNE, Proposed-Med, and Proposed-Mean) are applied to every session using 30-second epochs and all available channels. All 32 channels are included to ensure a fair comparison with BLINKER, which, by default, processes all channels provided to it and requires the user to select a representative blink channel. For consistency across methods, all pipelines are therefore supplied with the same 32 channels, and the channel yielding the highest $F_1$ score is selected for subsequent analysis. The per-session macro-$F_1$, precision, and recall are summarised in Table~\ref{tab:exp1_main}.

In the strategy comparison, both proposed configurations achieved higher macro-$F_1$ scores than BLINKER and MNE on Raja, Cao2018, and the pooled evaluation (Table~\ref{tab:exp1_main}). Proposed-Med achieved macro-$F_1$ scores of 0.8825, 0.8068, and 0.8403, respectively, whereas Proposed-Mean achieved 0.8742, 0.8030, and 0.8345. The corresponding scores were 0.7654, 0.6940, and 0.7256 for BLINKER and 0.6664, 0.5627, and 0.6086 for MNE. This performance gap is reflected in the precision--recall profiles shown in Figure~\ref{fig:condition_prf}. BLINKER achieved the highest recall among all conditions on Raja, Cao2018, and the pooled evaluation, with values of 0.9518, 0.9715, and 0.9627, respectively, but also exhibited the lowest precision, at 0.6794, 0.5696, and 0.6182. Wilcoxon signed-rank tests on session-level $F_1$, with Bonferroni correction across the six pairwise comparisons, showed significant differences between each proposed configuration and each baseline. Proposed-Med versus BLINKER yielded $p=1.32\times10^{-11}$, and Proposed-Med versus MNE yielded $p=6.51\times10^{-9}$. Similarly, Proposed-Mean versus BLINKER yielded $p=4.97\times10^{-10}$, while Proposed-Mean versus MNE yielded $p=9.94\times10^{-9}$. In contrast, neither within-pair comparison was significant: BLINKER versus MNE yielded $p=0.812$, whereas Proposed-Mean versus Proposed-Med yielded $p=1$. Furthermore, no channel group in either dataset produced an inversion in which a baseline matched or exceeded the performance of Proposed-Med (Table~\ref{tab:exp2_inversions}).

